\section{Phase geometry and mixed conditioning}\label{inf:sec-conditioning}

Fix the compact split interval of Section~\ref{inf:sec-widened}
and an integer split $a,b$ in it, with $a+b=2p\in\mathbb N$.
In this section $R\in[2p-1,2p+4]$ is the continued primary Neumann
zero, $0<\tau\le\log p$, and $|p-p_0|=O(\log p/p)$.
All assertions are for sufficiently large $p$, with constants
independent of $\tau$ in that range. Fix
$1<q_{\mathrm{gap}}<2$ and $0<\xi_{\mathrm{inv}}<2-q_{\mathrm{gap}}$.
Assume that the actual quartic domain has invariant gap at least
\[
 \gamma=\kappa_{\mathrm{gap}}\tau p^{-q_{\mathrm{gap}}},\qquad
 \kappa_{\mathrm{gap}}>0\text{ fixed}.
\]
The low-state structure is that of
Lemma~\ref{inf:wide-low-transfer} and the finite recurrence
\eqref{inf:finite-q-recurrence}.

\subsection{One radial pole at each angular index}

\begin{lemma}[Radial projection and complement]\label{inf:radial-poles}
For $0\le j\le\lfloor2p\rfloor$, choose the positive zero $z_j$
of $J_{\nu_j}$ closest in frequency to $R$, choosing the smaller
root in a tie. Let $\Pi_{\mathrm{rad}}$ retain its unit radial vector
and let $Q_{\mathrm{rad}}=I-\Pi_{\mathrm{rad}}$. These are contractions
in $\Mmix$ and $\Hc$. Put $d_j^{\mathrm{pole}}=z_j^2-R^2$. Then
\begin{equation}\label{inf:radial-pole-bounds}
 \nrm{K_0\{Q_{\mathrm{rad}}(H_0-R^2)Q_{\mathrm{rad}}\}^{-1}
                        Q_{\mathrm{rad}}}\le Cp,
 \qquad \nrm{K_0\Pi_{\mathrm{rad}}}\le Cp^2
\end{equation}
in both spaces and their angular-diagonal matrix norms.
On the retained space the mixed norm is exactly $\ell^1(2^j)$;
the Hilbert metric is $\diag(\eta_j^2)$.
\end{lemma}
\begin{proof}
The projections are radial orthogonal projections in each angular
block. Lemma~\ref{inf:mixed-algebra} proves the contraction statements.
Spacing greater than $\pi$ implies that every nonretained zero in a
retained block is at distance at least $\pi/2$ from $R$.
For $j>2p$ one has $\nu_j>5p-1$, so the centrifugal lower bound
puts its entire spectrum above $(5p-1)^2$. The scalar spectral
ratio $(z^2+R^2)/|z^2-R^2|$ is therefore at most $Cp$ on the
complement.

The radial sine trial on $(1/2,1)$ in
\eqref{inf:radial-schrodinger} bounds the first eigenvalue by
$4\nu_j^2+4\pi^2$. Its first zero is thus at most $2\nu_j+7$.
If $R\le2\nu_j$, comparison with that first zero gives
$z_j\le\max\{2R,2\nu_j+7\}=O(p)$.
If $R>2\nu_j$, the differential equation for
$\sqrt xJ_{\nu_j}(x)$ has frequency square
$1-(\nu_j^2-1/4)/x^2\ge3/4$ for $x\ge R$.
Comparison with a sine on an interval of length
$2\pi/\sqrt3$ gives a zero within a bounded distance beyond $R$.
For a direct proof, put $\kappa_{\mathrm{cmp}}=\sqrt3/2$ and
$w(x)=\sin(\kappa_{\mathrm{cmp}}(x-a))$ on
$[a,a+\pi/\kappa_{\mathrm{cmp}}]$. If a solution $v$ of
$v''+Q(x)v=0$, $Q\ge\kappa_{\mathrm{cmp}}^2$, had no zero there,
choose its sign positive. Then
$(v'w-vw')'=-(Q-\kappa_{\mathrm{cmp}}^2)vw\le0$,
whereas the difference of the endpoint values is
$\kappa_{\mathrm{cmp}}(v(a)+v(a+\pi/\kappa_{\mathrm{cmp}}))>0$.
This contradiction proves the comparison.
Again the closest zero has size $Cp$. Thus its $K_0$ eigenvalue is
at most $Cp^2$. The norm identifications follow from the radial
normalization and \eqref{inf:true-H-norm}.
\end{proof}

\begin{lemma}[Linear quartic reduction]\label{inf:centre-form-comparison}
Let $H_{\mathrm q}=z_{\mathrm q}S_2$ be the harmonic extension of
$f_{\mathrm q}$, so $|z_{\mathrm q}|\le C\tau p^{-2}$, and put
$L_{\mathrm q}=H_0+V_1(H_{\mathrm q})-R^2$.
The radial complementary inverse exists and its exact pole Schur
matrix is
\begin{equation}\label{inf:pole-schur-bands}
 \mathcal S_{\mathrm{pole}}
 =\diag(d_j^{\mathrm{pole}})+\mathcal B_{\mathrm{pole},2}
       +p^{-1}\mathcal B_{\mathrm{pole},4}
       +\mathsf R_{\mathrm{pole},6},
\end{equation}
where the first two bands have bandwidths two and four and
\begin{equation}\label{inf:pole-band-bounds}
 \nrm{\mathcal B_{\mathrm{pole},2}}\le C\tau,\quad
 \nrm{\mathcal B_{\mathrm{pole},4}}\le C\tau^2,\quad
 \nrm{\mathsf R_{\mathrm{pole},6}}\le C\tau^3p^{-2}
\end{equation}
in the coefficient column norm and Hilbert row/column norms.
Define the second-order interaction truncation by
\[
 \mathcal S_{\mathrm{pole}}^{(0)}
 =\diag(d_j^{\mathrm{pole}})+\mathcal B_{\mathrm{pole},2}
                         +p^{-1}\mathcal B_{\mathrm{pole},4}.
\]
It deletes the entire remainder, which contains all paths with
at least three interactions and need not have bandwidth six.
The exact matrix and this bandwidth-four truncation have Hilbert
gap at least $c\gamma$.
\end{lemma}
\begin{proof}
For this degree-four harmonic chart the first variation gives
\[
 V_1(H_{\mathrm q})\KD
 =2z_{\mathrm q}M_{S_2}
 +2z_{\mathrm q}\nabla S_2\cdot\nabla\Euler\KD
 +(2p+4)z_{\mathrm q}\nabla S_2\cdot\nabla\KD.
\]
The exact radial profiles in Lemma~\ref{inf:gradient-profiles}
factor each angular block into $c_{2,j}^i$ times a radial operator
of norm at most $C|z_{\mathrm q}|$.
Lemmas~\ref{inf:mixed-multipliers} and \ref{inf:mixed-Poisson}
bound its relative graph norm by $C\tau p^{-2}$ in both matrix
algebras. By \eqref{inf:radial-pole-bounds}, the radial complementary
Neumann ratio is $C\tau/p$. Direct projection has size $C\tau$;
one complementary inverse gives size $C\tau^2/p$ and bandwidth four.
Each further inverse/perturbation factor contributes $C\tau/p$.
Summing the convergent series proves \eqref{inf:pole-band-bounds},
including all complementary radial states.

For the spectral comparison the exact unitary form is
\[
 \int A_q(\nabla u-u\nabla\log m)\cdot
                  \overline{(\nabla u-u\nabla\log m)}.
\]
Its linear density term integrates to
$\int\Delta l_{\mathrm{lin}}|u|^2=0$.
The linear form has principal coefficient
$(1-2H_{\mathrm q})I-y\otimes\nabla H_{\mathrm q}
-\nabla H_{\mathrm q}\otimes y$ and is uniformly comparable
to the ball form. The remaining form error
is bounded by
$C|z_{\mathrm q}|^2(\nrm{\nabla u}_2^2+p^2\nrm u_2^2)$.
The ball Poincar\'e bound absorbs the second term into the first.
Indexwise min--max therefore changes eigenvalues near $R^2$ by at most
$C\tau^2p^{-2}=o(\gamma)$, since
$\tau p^{q_{\mathrm{gap}}-2}\to0$ uniformly here.
Thus $L_{\mathrm q}$ has at least half the actual quartic gap.
For a pole vector reconstruct the radial complementary component;
the full image is exactly its Schur image and its Hilbert norm
is at least that of the pole vector. This gives the Schur gap as
in Proposition~\ref{inf:inverse-Schur}.
Deleting $\mathsf R_{\mathrm{pole},6}$ changes that gap by
$C\tau^3p^{-2}=o(\gamma)$, proving the last assertion.
\end{proof}

\subsection{The two phase estimates}

Use the corrected phase
\begin{equation}\label{inf:corrected-phase}
 \varphi_{\mathrm{Bes}}(\nu,x)=\Xi(\nu,x)-\arctan b(\nu,x)
\end{equation}
with the functions in \eqref{inf:phase-def}. Chebyshev polynomials
of the second kind are denoted by $U_k^{\mathrm{Ch}}$, to distinguish
them from the Debye polynomials $U_k$.

\begin{lemma}[Two-sided phase and energy comparison]
\label{inf:phase-eigenvalue-two-sided}
Fix $0<c_*<1$. For $p-1\le\nu\le c_*R$, let $z$ be the zero of
$J_\nu$ nearest to $R$. If
$d_{\varphi}=\dist(\varphi_{\mathrm{Bes}}(\nu,R),\pi(\Z+1/2))$,
then for all sufficiently large $p$,
\begin{equation}\label{inf:phase-energy-comparison}
 c p(d_{\varphi}-Cp^{-2})
 \le |z^2-R^2|\le C p(d_{\varphi}+Cp^{-2}).
\end{equation}
The constants depend only on $c_*$.
\end{lemma}
\begin{proof}
Put $c_1=\sqrt{1-c_*^2}/2$ and choose a fixed interval
$|x-R|\le X_*$, $X_*=2\pi/c_1$. There $\kappa\ge cp$.
Lemma~\ref{inf:debye-finite} at order one gives
\[
 \frac{J_\nu(x)}{\sqrt{2/\pi}\kappa^{-1/2}}
 =\sqrt{1+b^2}\cos\varphi_{\mathrm{Bes}}(\nu,x)+e(x),
 \qquad |e(x)|\le C p^{-2}.
\]
The explicit phase derivatives give
$c_1\le\partial_x\varphi_{\mathrm{Bes}}\le1$ after increasing
$p$. At a zero the distance of its phase from $\pi(\Z+1/2)$ is
at most $Cp^{-2}$. Conversely, on either side of a lattice phase at
phase distance $2Cp^{-2}$, enlarge $C$ to dominate the error above.
The two values have opposite signs, so the intermediate value
theorem supplies a zero between them. Every lattice phase whose
distance from both endpoints of the phase interval is greater
than $2Cp^{-2}$ has such a zero inside the argument interval.
Since $\partial_x\varphi_{\mathrm{Bes}}\ge c_1$ and
$X_*=2\pi/c_1$, the phase interval contains
$[\varphi_{\mathrm{Bes}}(\nu,R)-2\pi,
\varphi_{\mathrm{Bes}}(\nu,R)+2\pi]$.
The lattice phase nearest to $\varphi_{\mathrm{Bes}}(\nu,R)$ is
at distance at most $\pi/2$, so it has a fixed positive endpoint
margin and satisfies this condition for all sufficiently large $p$.

Choose the lattice phase nearest to the phase at $R$ and a zero
$z'$ associated to it. The derivative lower bound gives
$|z'-R|\le(d_{\varphi}+Cp^{-2})/c_1$. Thus the nearest frequency
zero $z$ also lies in the chosen interval and obeys that upper bound.
For its lower bound, its phase is within $Cp^{-2}$ of some lattice
phase, whence
$d_{\varphi}\le|\varphi_{\mathrm{Bes}}(\nu,z)
-\varphi_{\mathrm{Bes}}(\nu,R)|+Cp^{-2}\le|z-R|+Cp^{-2}$.
Finally $cp\le z+R\le Cp$ on this interval. Multiplication proves
\eqref{inf:phase-energy-comparison}. This proof uses the order-one
finite expansion for all $x/\nu>1$ bounded away from one, including
the ratios below $3/2$ required for the high critical regions.
\end{proof}

\begin{lemma}[Finite-shift resonance through Lommel polynomials]
\label{inf:lommel-resonance}
Fix $0<\alpha_{\mathrm{loc}}<1$. If retained indices $j,j+d$,
$1\le d\le4$, both satisfy
$|d_j^{\mathrm{pole}}|,|d_{j+d}^{\mathrm{pole}}|\le p^{1-\alpha_{\mathrm{loc}}}$,
then
\[
 |U_{2d-1}^{\mathrm{Ch}}(\nu_j/R)|\le Cp^{-\alpha_{\mathrm{loc}}}.
\]
The ratio $\nu_j/R$ is within $Cp^{-\alpha_{\mathrm{loc}}}$ of
\begin{equation}\label{inf:critical-ratios}
 \mathcal T_{\mathrm{crit}}=
 \{1/2,1/\sqrt2,\sqrt3/2,\cos(\pi/8)\}.
\end{equation}
In particular $d=1$ is impossible for large $p$.
\end{lemma}
\begin{proof}
Write $\nu=\nu_j$, $\mu=\nu+2d$, and let $z,z'$ be the two retained
zeros. Their distances from $R$ are at most $Cp^{-\alpha_{\mathrm{loc}}}$.
At $J_\nu(z)=0$, one has $J_{\nu+1}(z)\ne0$ by ODE uniqueness.
Dividing the Bessel recurrence by this nonzero value gives the
finite Lommel recurrence, with initial values zero and one at
shifts zero and one. This identity holds for every real order.
For the finitely many shifts at most nine its recurrence coefficients
are $2\nu/z+O(p^{-1})$. The recurrence
$U_{k+1}^{\mathrm{Ch}}(t)=2tU_k^{\mathrm{Ch}}(t)-U_{k-1}^{\mathrm{Ch}}(t)$
therefore gives, by induction,
\[
 \frac{J_{\nu+m}(z)}{J_{\nu+1}(z)}
 =U_{m-1}^{\mathrm{Ch}}(\nu/z)+O(p^{-1}),\qquad 1\le m\le9.
\]
The ratios and the derivative ratio of $J_\mu$ are bounded by an
absolute constant there, using the derivative recurrence.
On the interval joining $z,z'$, the coefficients of the system for
$(J_\mu,J'_{\mu})$ are uniformly bounded: $x\asymp p$ and
$\mu<z'\le R+Cp^{-\alpha_{\mathrm{loc}}}$.
Its integral equation and Gronwall's inequality give
$\sup|J_\mu'|\le C|J_{\nu+1}(z)|$ on that interval.
Since $J_\mu(z')=0$, integration proves
$|J_\mu(z)|\le Cp^{-\alpha_{\mathrm{loc}}}|J_{\nu+1}(z)|$.
Divide and use the preceding recurrence; replacing $z$ by $R$
changes its polynomial argument by $O(p^{-1-\alpha_{\mathrm{loc}}})$.
This proves the bound, including near the turning region.

The zeros of $U_{2d-1}^{\mathrm{Ch}}$ are the simple numbers
$\cos(h\pi/(2d))$, $1\le h<2d$; this follows from
$U_k^{\mathrm{Ch}}(\cos\beta)=\sin((k+1)\beta)/\sin\beta$.
Here $\nu/R\ge1/2-o(1)$ and is at most $1+o(1)$.
The roots in this range, for $d=1,2,3,4$, are respectively
none, $1/\sqrt2$, $1/2,\sqrt3/2$, and
$1/\sqrt2,\cos(\pi/8)$. Simplicity and
$U_{2d-1}^{\mathrm{Ch}}(1)=2d$ prove the remaining assertions.
\end{proof}

\subsection{Core geometry and Green propagation}

Choose $\alpha_{\mathrm{loc}}>0$ with
$q_{\mathrm{gap}}+\xi_{\mathrm{inv}}<2-2\alpha_{\mathrm{loc}}$.
Set $W_{\mathrm{cond}}=C_{\mathrm{cond}}(1+\tau)$, with a fixed
sufficiently large $C_{\mathrm{cond}}$, and
\[
 \mathcal I_{\mathrm{win}}=\{j:|d_j^{\mathrm{pole}}|\le W_{\mathrm{cond}}\}.
\]
The angular complement of this set in the bandwidth-four matrix
$\mathcal S_{\mathrm{pole}}^{(0)}$ has every principal inverse
bounded by $2/W_{\mathrm{cond}}$. Indeed its band perturbation has
norm $C\tau$ and division by the unperturbed diagonal has norm at most
$1/W_{\mathrm{cond}}$, also after any coordinate restriction.
Let $\mathcal F_{\mathrm{win}}$ be the Schur complement of
$\mathcal S_{\mathrm{pole}}^{(0)}$ onto $\mathcal I_{\mathrm{win}}$.
It is self-adjoint for $\diag(\eta_j^2)$ and has Hilbert gap $c\gamma$
by Lemma~\ref{inf:centre-form-comparison} and angular complementary
reconstruction. The remainder $\mathsf R_{\mathrm{pole},6}$ is
restored after the window inverse in the proof of
Theorem~\ref{inf:mixed-conditioning}.

\begin{lemma}[Critical strips and the true minima]\label{inf:critical-strips}
For $c_{\mathrm{cr}}\in\mathcal T_{\mathrm{crit}}$ let
\[
 \mathcal N_{c_{\mathrm{cr}}}
 =\{j\in\R:|p+2j-1-Rc_{\mathrm{cr}}|
                       \le2p^{1-\alpha_{\mathrm{loc}}/2}\}.
\]
These neighbourhoods are disjoint for large $p$.
Write $c_{\mathrm{cr}}=\cos(\ell_{c_{\mathrm{cr}}}\pi/(2d_{c_{\mathrm{cr}}}))$
with reduced pairs
$(\ell_{c_{\mathrm{cr}}},d_{c_{\mathrm{cr}}})=(2,3),(1,2),(1,3),(1,4)$,
respectively, and put
\[
 f_{c_{\mathrm{cr}}}(j)=\varphi_{\mathrm{Bes}}(p+2j-1,R)
                      +\ell_{c_{\mathrm{cr}}}\pi j/d_{c_{\mathrm{cr}}}.
\]
On $\mathcal N_{c_{\mathrm{cr}}}$ its second derivative is between
$c_{\mathrm{curv}}/p$ and $C_{\mathrm{curv}}/p$, where
$c_{\mathrm{curv}},C_{\mathrm{curv}}>0$ are fixed, and
$|f'_{c_{\mathrm{cr}}}|\le Cp^{-\alpha_{\mathrm{loc}}/2}$.
Its true minimum is attained at
\begin{equation}\label{inf:true-phase-minimum}
 j_{c_{\mathrm{cr}}}^*=\{Rc_{\mathrm{cr}}-(p-1)\}/2+O(p^{-1}).
\end{equation}
Connect successive window vertices in each high neighbourhood
$c_{\mathrm{cr}}>1/2$ if their distance is at most
$L_{\mathrm{core}}(p)=K_{\mathrm{core}}\log p/\log\log p$.
Each resulting component has a common phase level and residue and
diameter at most
\begin{equation}\label{inf:sharp-core-diameter}
 L_{\mathrm{core}}(p)+C\sqrt{W_{\mathrm{cond}}+1}.
\end{equation}
Its indices satisfy $j\ge c_2p$ for a fixed $c_2>0$.
The low core consists of window indices $j\le C\sqrt{W_{\mathrm{cond}}+1}$,
all congruent to two modulo three. Every other low-neighbourhood
window vertex has $j\ge c\sqrt p$ and is isolated at scale
$L_{\mathrm{core}}(p)$.
For the nonseed indices of the low core one also has the uniform formula
\begin{equation}\label{inf:low-core-offsets}
 \begin{gathered}
 d_j^{\mathrm{pole}}=-\frac43(2j-1)(2j-4)
             +O\bigl((1+\log p)^2/p\bigr),\\
 j\equiv2\pmod3,\qquad 5\le j\le C\sqrt{W_{\mathrm{cond}}+1}.
 \end{gathered}
\end{equation}
For $j=5,8,11$, these leading offsets are $-72,-240,-504$, respectively.
\end{lemma}
\begin{proof}
The derivatives in the proof of Lemma~\ref{inf:phase-strip} give
$f_{c_{\mathrm{cr}}}''=4/\kappa+O(p^{-3})$. At
$j_{c_{\mathrm{cr}}}=\{Rc_{\mathrm{cr}}-(p-1)\}/2$ the derivative of the leading phase cancels
the added linear term, while the corrected derivative is
$-2b_\nu/(1+b^2)<0$, of size $O(p^{-2})$.
Its strictly positive second derivative
therefore puts its unique minimum at distance $O(p^{-1})$, proving
\eqref{inf:true-phase-minimum}.

By Lemma~\ref{inf:phase-eigenvalue-two-sided}, window vertices lie
within $C(W_{\mathrm{cond}}+1)/p$ of a level of the lattice
$\pi(\Z+1/2)+\ell_{c_{\mathrm{cr}}}\pi j/d_{c_{\mathrm{cr}}}$.
Levels are spaced by $\pi/d_{c_{\mathrm{cr}}}$.
An edge of length at most $L_{\mathrm{core}}$ changes $f_{c_{\mathrm{cr}}}$ by
$o(1)$, so its endpoints have the same level; reduction of
$(\ell_{c_{\mathrm{cr}}},d_{c_{\mathrm{cr}}})$ makes their residues equal.
Every component thus lies in one strip
$|f_{c_{\mathrm{cr}}}-g|\le C(W_{\mathrm{cond}}+1)/p$.
If $f''\ge m>0$, either monotone side of $|f-g|\le\delta$
has length at most $2\sqrt{\delta/m}$: its variation over length
$h$ is at least $mh^2/2$ and at most $2\delta$.
A connected chain either stays on one side or has an edge of
length at most $L_{\mathrm{core}}$ joining the two sides.
This proves \eqref{inf:sharp-core-diameter}. High neighbourhoods
have $j\asymp p$ by their definitions.

For $c_{\mathrm{cr}}=1/2$ the seed $j=2$ anchors the first level, since
$d_2^{\mathrm{pole}}=-\theta/(2p)+O(\theta^2p^{-3})$.
For $j\le c\sqrt p$ the phase cannot reach another level.
The curvature bound then gives $j\le C\sqrt{W_{\mathrm{cond}}+1}$,
and the first-level residue is $j\equiv2\pmod3$.
Beyond this first group, a window vertex has $j\ge c\sqrt p$.
There $f'_{c_{\mathrm{cr}}}\ge cj/p$; moving by an integer between one and
$L_{\mathrm{core}}$ changes its level by more than
$2C(W_{\mathrm{cond}}+1)/p$ and less than
$\pi/3-2C(W_{\mathrm{cond}}+1)/p$.
Hence such vertices are isolated at that scale.
To obtain \eqref{inf:low-core-offsets}, apply
Lemma~\ref{inf:low-recurrence} at the associated primary crossing
$p_0$. For $j\equiv2\pmod3$, $j\ge5$, its value/normal ratio is
$c_j/p_0+O((j+1)^4/p_0^2)$: the numerator has the stated
first-order recurrence coefficient, and the normal denominator
is its nonzero limiting sign plus $O((j+1)^2/p_0)$.
The collar ODE is uniformly bounded for
$j=O(\sqrt{\log p})$; its nearby zero therefore has squared
offset $-4c_j+O((j+1)^4/p_0)$.
Continue this zero and the Neumann zero to $p$.
Their frequency difference stays $O(\log p/p)$ by the global
order derivative bound, so spacing identifies the continued zero
with the retained pole. Repeating the refined derivative comparison
in the proof of Lemma~\ref{inf:wide-low-transfer} gives
$|z'(t)-R'(t)|\le C(j+1)/p+C(\log p)/p^2$.
The interval has length $O(\log p/p)$; the change of squared
offset is consequently at most
$C(\log p)(j+1)/p+C(\log p)^2/p^2$.
Since $j\le C\sqrt{W_{\mathrm{cond}}+1}=O(\sqrt{\log p})$,
these errors give \eqref{inf:low-core-offsets}.
\end{proof}

\begin{lemma}[Bin-elimination Green bound]\label{inf:bin-green}
Let a matrix on a finite subset of integer sites be
$\diag(d_j)+B$, where $B$ has bandwidth four and norm at
most $b_*\ge0$ in a weighted coordinate $\ell^1$ norm or a diagonal
Hilbert norm. Suppose $W>0$, $|d_j|>W$, $b_*/W<1/4$, and all coordinate
compressions are contractions. Assume the same band bound on
principal restrictions. Divide the sites into consecutive bins of
width four. Every principal inverse has norm at most $2/W$.
If a run of $m$ intervening bins has all $|d_j|\ge L>0$, with
$L\ge4(b_*+2b_*^2/W)$, then its bin-to-bin Green propagation has
bound
\begin{equation}\label{inf:bin-green-bound}
 \frac{C}{W}(2b_*/L)^m.
\end{equation}
Without the stronger diagonals there is exponential decay in the
number of intervening bins. Empty bins disconnect the matrix.
\end{lemma}
\begin{proof}
Diagonal division and a Neumann series bound every principal inverse
by $1/(W-b_*)\le2/W$. The bin matrix is block tridiagonal.
Eliminate bins from the right. The inverse of each right Schur
block is the compression of the corresponding tail inverse, so
has norm at most $2/W$. The diagonal block at a strong bin differs
from the restriction of $\diag(d_j)$ by at most $b_*+2b_*^2/W$.
Another diagonal Neumann series therefore bounds its inverse by
$2/L$. On a homogeneous row, backward elimination gives exactly
$u_b=-G_bT_{b,b-1}u_{b-1}$, with $\nrm{T_{b,b-1}}\le b_*$.
Iteration across the strong run gives $(2b_*/L)^m$.
For a source bin, eliminate on both sides; its resulting solution
operator is a compression of the full inverse and costs $2/W$.
This proves \eqref{inf:bin-green-bound}. Replacing $L$ by the
ordinary Schur bound gives ratio $2b_*/W<1/2$ on other bins.
The argument is valid separately in either specified norm; no
change of coordinates occurs inside the propagation estimate.
\end{proof}

\begin{lemma}[Valuations of low angular paths]\label{inf:path-valuation}
Uniformly for $j\le C\sqrt{1+\tau}+4$, the quartic angular
coefficients satisfy
\begin{equation}\label{inf:low-coefficient-valuations}
 \begin{aligned}
 c_{2,j}^{j+2}&=O(1),&
 c_{2,j}^{j+1},c_{2,j}^j&=O((j+1)/p),\\
 c_{2,j}^{j-1},c_{2,j}^{j-2}&=O((j+1)^2/p^2).
 \end{aligned}
\end{equation}
Every path of quartic interactions and angular-diagonal radial
resolvents carries the product of its angular coefficients.
This factorization applies to the feedback band
$p^{-1}\mathcal B_{\mathrm{pole},4}$ and to every further complementary path.
A two-interaction path starting at a low index $j$ in the stated
range has an extra factor $C(j+1)/p$ in its angular coefficient
product in each of the two cases of net displacement $+3$ and $-3$.
\end{lemma}
\begin{proof}
Put
\[
 \mathfrak u_j=\frac{\mathsf H_{j+1}(1)}{\mathsf H_j(1)},\qquad
 \mathfrak v_j=\mathfrak a_j^2
              \frac{\mathsf H_{j-1}(1)}{\mathsf H_j(1)},\qquad
 a_{\mathrm{quad}}=\frac{2(\ab+1)}{p+2},
\]
with $\mathfrak v_0=0$. The normalized recurrence is
$xG_j=\mathfrak u_jG_{j+1}+\mathfrak b_jG_j+\mathfrak v_jG_{j-1}$.
Two applications give the exact coefficients
\begin{align*}
 \mathsf H_2(1)c_{2,j}^{j+2}&=\mathfrak u_j\mathfrak u_{j+1},\\
 \mathsf H_2(1)c_{2,j}^{j-2}&=\mathfrak v_j\mathfrak v_{j-1},\\
 \mathsf H_2(1)c_{2,j}^{j+1}
 &=\mathfrak u_j(\mathfrak b_j+\mathfrak b_{j+1}-a_{\mathrm{quad}}),\\
 \mathsf H_2(1)c_{2,j}^{j-1}
 &=\mathfrak v_j(\mathfrak b_j+\mathfrak b_{j-1}-a_{\mathrm{quad}}),\\
 \mathsf H_2(1)c_{2,j}^{j}
 &=\mathfrak u_j\mathfrak v_{j+1}+\mathfrak v_j\mathfrak u_{j-1}
                                      +\mathsf H_2(\mathfrak b_j).
\end{align*}
Terms outside the nonnegative index range are zero.
The formulas \eqref{inf:recurrence-a}--\eqref{inf:recurrence-b}
and \eqref{inf:angular-jacobi} give explicitly
\begin{align*}
 \mathfrak u_j&=\frac{(p(1-r)+j)(p+j-1)}{(p+2j)(p+2j-1)},\\
 \mathfrak v_j&=\frac{j(pr+j-1)}{(p+2j-1)(p+2j-2)},\\
 \mathfrak b_j-r&=\frac{2j(p+j-1)(1-2r)}{(p+2j)(p+2j-2)},
 \qquad a_{\mathrm{quad}}-2r=\frac{2(1-2r)}{p+2}.
\end{align*}
Their denominators are comparable to the displayed powers of $p$
uniformly for $j\le C\sqrt{1+\tau}+4$. In particular
$\mathfrak u_j=O(1)$, $\mathfrak v_j=O((j+1)/p)$,
$\mathfrak b_j-r=O((j+1)/p)$, and $a_{\mathrm{quad}}-2r=O(p^{-1})$.
The endpoint ratios are bounded above and below on this range,
and $\mathsf H_2(1)\ge c_{\mathrm{sp}}>0$.
Writing \eqref{inf:H2} in powers of $x-r$ gives
$\mathsf H_2(\mathfrak b_j)=O((j+1)^2/p^2+p^{-1})$.
Since $j=O(\sqrt{\log p})$, every displayed coefficient has exactly
the bound in \eqref{inf:low-coefficient-valuations}.

The exact block formula in the proof of
Lemma~\ref{inf:centre-form-comparison} factors one interaction as
$c_{2,j}^i$ times its radial operator. An angular-diagonal radial
resolvent preserves this scalar factor. Induction on the number
of interactions therefore gives the product on any path, followed
by the product of its bounded radial norms. The Hilbert endpoint
ratio $\eta_i/\eta_j$ telescopes along the path.
Two steps of lengths at most two with net displacement three
must include a one-step edge. Equation~\eqref{inf:low-coefficient-valuations}
gives the asserted extra factor. Positive coefficient sums and
bounded bandwidth convert these bounds to column estimates,
without a factor for the number of low indices.
\end{proof}

\begin{lemma}[The seven off-core block types]\label{inf:off-core-table}
Take the high components of Lemma~\ref{inf:critical-strips}, its
low core, and all remaining window vertices as singleton cores.
For sufficiently large fixed $K_{\mathrm{core}}$, in both the
coefficient operator norm and the Hilbert operator norm,
\begin{equation}\label{inf:global-off-core}
 \nrm{\operatorname{Off}_{\mathrm{cores}}\mathcal F_{\mathrm{win}}}
 \le C\tau(1+\tau)^2p^{-2+2\alpha_{\mathrm{loc}}}.
\end{equation}
The low-core internal off-diagonal part has the stronger bound
$C\tau(1+\tau)^2p^{-2}$ in both norms.
\end{lemma}
\begin{proof}
Call an outside singleton one outside all $\mathcal N_{c_{\mathrm{cr}}}$, and a low
singleton one in $\mathcal N_{1/2}$ beyond its low core.
The following seven types exhaust all pairs of distinct cores;
the last row takes precedence when the neighbourhoods differ.
\begin{center}\small\normalfont
\begin{tabular}{@{}>{\raggedright\arraybackslash}p{.37\textwidth}>{\raggedright\arraybackslash}p{.56\textwidth}@{}}
\toprule Block type & Estimate and mechanism\\\midrule
Outside--outside & $C\tau(1+\tau)^2p^{-2+2\alpha_{\mathrm{loc}}}$;
two attenuated endpoints.\\
Outside--low singleton & Same bound; the low endpoint has the
stronger attenuation $C\tau/p$.\\
Low singleton--low singleton & Same bound; two primary edges
cannot join distinct such vertices.\\
Outside--high component & Faster than any chosen power,
including pairs crossing a critical-neighbourhood boundary.\\
High component--high component in one neighbourhood &
Convex phase barrier; $C\tau p^{-B_*}$ for any prescribed fixed $B_*$.\\
Low core--other core in its neighbourhood or outside &
Generic Green decay across distance at least $c\sqrt p$.\\
Cores in different critical neighbourhoods &
Generic Green decay across distance at least $cp$.\\
\bottomrule
\end{tabular}
\end{center}

Here are the operator estimates behind the table. An outside vertex
has every complementary neighbour within four sites at energy
distance greater than $p^{1-\alpha_{\mathrm{loc}}}$, by
Lemma~\ref{inf:lommel-resonance}; otherwise it would lie in the
smaller critical neighbourhood of radius $Cp^{1-\alpha_{\mathrm{loc}}}$.
Thus either endpoint map, with its adjacent diagonal inverse,
has norm $C\tau p^{-1+\alpha_{\mathrm{loc}}}$.
At a low singleton, for $h=\pm1,\pm2$ the phase change to $j+h$
is $-2h\pi/3+o(1)$, because $j/p\to0$ in the low neighbourhood
and $\partial_\nu\varphi_{\mathrm{Bes}}=-\arccos(\nu/R)+O(p^{-2})$.
Thus primary neighbours have phase separated from the zero lattice
by a fixed amount and denominators at least $cp$.
Feedback edges
already have size $C\tau^2/p$, and their denominator is at least
$W_{\mathrm{cond}}$. They give the endpoint bound $C\tau/p$.
All these maps have fixed bandwidth, so the estimates hold on
whole columns and, in Hilbert coordinates, on rows as well.

No direct band-four entry connects two distinct singleton cores.
Nor can two primary edges connect them: the endpoints would be
two window vertices within four sites and hence in the same
critical component. A length-two path between distinct singletons
therefore includes a feedback edge. With one such edge put the
diagonal inverse at the primary endpoint, giving
$C\tau^3p^{-2+\alpha_{\mathrm{loc}}}$; with two use
$C\tau^4/(p^2W_{\mathrm{cond}})$.
For all paths of length at least three, the resolvent identity
places a diagonal inverse at each endpoint and a bounded
Neumann inverse between them. Their global bound is
$C\tau^3p^{-2+2\alpha_{\mathrm{loc}}}$.
This proves the first three rows as operator norms, with no
sum over the total number of vertices.

For two distinct high components in one neighbourhood, let their
nearest endpoints be $u<v$, with $D=v-u>L_{\mathrm{core}}$, and set
\[
 h=\pi/d_{c_{\mathrm{cr}}},\qquad
 \delta=C(W_{\mathrm{cond}}+1)/p.
\]
Here $h$ is the spacing of the union of the phase levels of all
integer residues. If their phase levels differ,
there is an intervening phase interval a fixed distance from all
levels. Since $|f'_{c_{\mathrm{cr}}}|\le Cp^{-\alpha_{\mathrm{loc}}/2}$,
this supplies at least $cp^{\alpha_{\mathrm{loc}}/2}$ consecutive
sites with denominator $cp$ by
\eqref{inf:phase-energy-comparison}. If both endpoints have level
$g$, define their excursion by
$\mathfrak e=g-\min_{[u,v]}f_{c_{\mathrm{cr}}}$.
The convexity inequality gives
$g-f_{c_{\mathrm{cr}}}(x)\ge cD^2/p-\delta$
on $[u+D/4,v-D/4]$.
If $\mathfrak e\ge h/4$, continuity supplies the phase interval
$[g-h/4,g-h/8]$, a fixed distance from every level;
use the preceding barrier on its preimage.
If $\mathfrak e<h/4$, the middle half has matching-residue
denominators at least $cD^2$, and all other residues have
denominators at least $cp$. This branch forces $D\le C\sqrt p$,
so, after decreasing its fixed constant, take $L=cD^2\le cp$.
Since $D>L_{\mathrm{core}}$ and
$L_{\mathrm{core}}^2/(1+\tau)\to\infty$, one has
$L\ge4(b_*+2b_*^2/W_{\mathrm{cond}})$ with $b_*=C\tau$.
The middle half contains at least $D/8-2$ complete bins of width
four. Their diagonals exceed $W_{\mathrm{cond}}$ for large $p$.
Lemma~\ref{inf:bin-green} bounds propagation
by $\exp\{-cD\log(D^2/[C(1+\tau)])\}$.
Uniformly for $\tau\le\log p$,
\[
 \log\frac{L_{\mathrm{core}}(p)^2}{C(1+\tau)}
 \ge\tfrac12\log\log p
\]
eventually. Choose $K_{\mathrm{core}}$ to give any fixed power
$p^{-B_*-2}$. There are only $O(p)$ retained vertices; summing
the already weighted entry bounds gives all coefficient columns
and Hilbert rows and columns with the desired power.
Each such Schur term has two exterior couplings. Together with
the source factor $1/W_{\mathrm{cond}}$ their size is at most
$C\tau^2/W_{\mathrm{cond}}\le C\tau$, giving the detuning factor
in the high-component rows of the table.

At a boundary of $\mathcal N_{c_{\mathrm{cr}}}$, $|f'_{c_{\mathrm{cr}}}|$ is comparable to
$p^{-\alpha_{\mathrm{loc}}/2}$. Two window points on opposite
sides of that boundary cannot have the same level at integer
distance between one and $cp^{\alpha_{\mathrm{loc}}/2}$:
their phase change exceeds $C(W_{\mathrm{cond}}+1)/p$.
Different levels require at least that distance as well.
Thus generic exponential Green decay proves the outside--high
row, including these boundary pairs. Distinct neighbourhoods are
separated by $cp$. The low core is separated from other vertices
of its neighbourhood by $c\sqrt p$. These facts prove the last
two rows.

It remains to bound the internal low-core off-diagonal part.
Its vertices have residue two modulo three, so there is no direct
primary edge between distinct vertices. The direct feedback of
displacement three has size
$C\tau^2(j+1)p^{-2}$ by Lemma~\ref{inf:path-valuation}.
A primary endpoint lands on a nonresonant index with denominator
$cp$, uniformly for $j\le C\sqrt{1+\tau}$.
Indeed $f_{1/2}(j)-f_{1/2}(2)=O((j+1)^2/p)=O(W_{\mathrm{cond}}/p)$,
while either of the other two residues has its zero lattice shifted
by $\pi/3$ from this level. Lemma~\ref{inf:phase-eigenvalue-two-sided}
therefore supplies that lower bound also for the neighbouring indices.
A feedback endpoint of displacement one, two or four also lands
on a nonresonant index; displacement three retains the extra
coefficient $(j+1)/p$ even if its denominator is only
$W_{\mathrm{cond}}$. Two-primary-edge paths of net displacement
three gain the same extra coefficient and a $cp$ denominator.
A zero-displacement step stays at its window index and is therefore
absent from an endpoint map into the angular complement.
Consequently each endpoint map has coefficient norm at most
$C\tau/p$. A two-edge term containing a feedback edge is bounded by
$C\tau^3p^{-2}+C\tau^4/(p^2W_{\mathrm{cond}})$, placing the inverse
at the primary endpoint when one is present.
The resolvent identity bounds the sum of paths with at least three
edges by $C\tau^3p^{-2}$, using the two endpoint bounds and a
bounded middle Neumann inverse. Bounded-band column sums therefore give
$C\tau(1+\tau)^2p^{-2}$ for the entire low-core off-diagonal matrix.
It is self-adjoint in a positive diagonal metric. Its Hilbert
operator norm equals its spectral radius, which is at most its
coefficient operator norm. This supplies the Hilbert bound without
any factor for the growing low-core cardinality.
\end{proof}

\begin{theorem}[Mixed conditioning and restoration]
\label{inf:mixed-conditioning}
Under the hypotheses of this section the linear quartic operator
has a two-sided mixed graph inverse with
\begin{equation}\label{inf:linear-mixed-inverse}
 \nrm{K_0L_{\mathrm q}^{-1}}_{\Mmix\to\Mmix}
 \le C p^{2+\xi_{\mathrm{inv}}}/\gamma.
\end{equation}
For a fixed order $N$, suppose its true centre has bounded degree,
strong size $C\tau p^{-2}$, and
\begin{equation}\label{inf:conditioning-centre-hypothesis}
 \nrm{f_N-f_{\mathrm q}}_{\mathrm{strong}}
 \le\tau p^{-4}\mathfrak E_N(p),\qquad
 \mathfrak E_N(p)=p^{o(1)}
\end{equation}
uniformly in the parameters. Then its actual scalar graph is a
bounded isomorphism on $\Cc$ and
\begin{equation}\label{inf:conditioning-final-inverse}
 \nrm{J_D^{-1}}_{\Cc\to\Cc}
 \le C_Np^{5/2+\xi_{\mathrm{inv}}}/\gamma.
\end{equation}
The threshold may depend on the compact split interval,
$N,q_{\mathrm{gap}},\xi_{\mathrm{inv}},\kappa_{\mathrm{gap}}$, and
the envelope in \eqref{inf:conditioning-centre-hypothesis}.
\end{theorem}
\begin{proof}
Delete cross-core entries of
$\mathcal F_{\mathrm{win}}$. Their Hilbert norm is $o(\gamma)$ by
\eqref{inf:global-off-core}, so every core retains Hilbert gap
$c\gamma$. In a high core, the exact ratio
\begin{equation}\label{inf:core-conditioning}
 \frac{\eta_{j+1}^2}{\eta_j^2}
 =\frac{2j+p-1}{2j+p+1}
   \frac{(pr+j)(j+1)}{(p(1-r)+j)(p-1+j)}
\end{equation}
is bounded above and below by fixed positive constants since
$j\ge c_2p$ and $j=O(p)$. In coordinates
$\widetilde u_j=2^ju_j$, the coefficient norm is unweighted
$\ell^1$ and the Hilbert weights are $\eta_j/2^j$.
Their consecutive ratios have the same boundedness property.
The diagonal similarity over the core diameter, followed by
$\ell^2$--$\ell^1$ comparison over its cardinality, therefore costs at most
\[
 \exp\{C L_{\mathrm{core}}(p)+C\sqrt{W_{\mathrm{cond}}+1}\}
 \sqrt{L_{\mathrm{core}}(p)+C\sqrt{W_{\mathrm{cond}}+1}+1}
 =p^{o(1)}.
\]
For a singleton the cost is one. On the low core, delete its
internal off-diagonal part in the Hilbert norm. Each diagonal
entry then has modulus at least $c\gamma$. Restoring that part
by a coefficient diagonal Neumann series gives inverse $C/\gamma$,
because $(1+\tau)^2p^{q_{\mathrm{gap}}-2}\to0$.

The direct sum of all core inverses has coefficient norm at most
$Cp^{\xi_{\mathrm{inv}}}/\gamma$, since the displayed subpolynomial
factor is at most $p^{\xi_{\mathrm{inv}}/2}$ for sufficiently large $p$. Its product with the off-core error is bounded by
$C(1+\tau)^2p^{q_{\mathrm{gap}}+\xi_{\mathrm{inv}}-2+2\alpha_{\mathrm{loc}}}=o(1)$.
Thus $\mathcal F_{\mathrm{win}}^{-1}$ has that same bound.
Recovering the angular complement costs bounded factors, since
$\tau/W_{\mathrm{cond}}$ is small. Restoring
$\mathsf R_{\mathrm{pole},6}$ costs
$C\tau^2p^{q_{\mathrm{gap}}+\xi_{\mathrm{inv}}-2}=o(1)$.
Hence the exact pole Schur inverse has the same coefficient bound.

The radial complementary reconstruction has reduced right side
of norm at most $(1+C\tau/p)\nrm f_{\Mmix}$.
The pole component costs $Cp^{\xi_{\mathrm{inv}}}\nrm f/\gamma$;
its $K_0$ norm costs a further $Cp^2$.
The complementary graph norm is bounded by
$Cp(\nrm f+C\tau\nrm{\Pi_{\mathrm{rad}}u})$.
This proves \eqref{inf:linear-mixed-inverse} on the entire graph
domain, with all radial states included. Lemma~\ref{inf:mixed-domain}
applies to the small relative perturbation $V_1(H_{\mathrm q})$;
the bounded scalar shift leaves its graph domain unchanged.
The reconstructions converge in that graph norm. The embeddings and
uniqueness identify this inverse with the Hilbert inverse on
mixed sources, and prove both inverse identities.

For the actual centre put $\Phi=\Phi_N$. Apply
Lemma~\ref{inf:quadratic-graph} with $H=\HP f_N$ and use
\eqref{inf:conditioning-centre-hypothesis}. They give
\[
 \nrm{(\Hu-H_0-V_1(H_{\mathrm q}))K_0^{-1}}
 \le C_N\tau p^{-4}(\mathfrak E_N(p)+\tau)=\tau p^{-4+o(1)}.
\]
Multiply by the mixed graph inverse before any radial coefficient
conversion. Its product is
$p^{q_{\mathrm{gap}}+\xi_{\mathrm{inv}}-2+o(1)}=o(1)$.
A Neumann series restores the true centre on $\Mmix$.
On $\Hc$, its linear graph inverse satisfies
$\nrm{K_0L_{\mathrm q}^{-1}}\le Cp^2/\gamma$.
Indeed, for $L_{\mathrm q}u=f$ write
$K_0u=f+2R^2u-V_1(H_{\mathrm q})u$, use the Hilbert gap
$\nrm u\le C\gamma^{-1}\nrm f$, and absorb
$\nrm{V_1(H_{\mathrm q})K_0^{-1}}<1/2$.
The same centre perturbation therefore has Hilbert product
$p^{q_{\mathrm{gap}}-2+o(1)}=o(1)$.
The restored Hilbert and mixed inverses agree on mixed sources
by the embedding and Hilbert uniqueness.
The density and graph-conjugation proof of
\eqref{inf:density-graph-bounds} uses only $p\delta$ small;
its binomial bound is
$\exp(Cp\delta)=\exp(O(\tau/p))$, hence bounded.
The scalar return therefore preserves the mixed inverse exponent.
The scalar estimates in the proof of Theorem~\ref{inf:relative-graph}
bound the sum of the $\Cc$ and $\Mmix$ relative norms by
$C(\delta+p\delta^2)\le C\tau p^{-2}$.
Thus \eqref{inf:head-tail-hypotheses} holds for large $p$.
The finite-head/infinite-tail proof of Lemma~\ref{inf:head-tail}
now costs exactly
$\sqrt p$. Its tail is solved in $\Cc$ and its head is converted
from the already constructed mixed solution. This proves
coefficient membership, both inverse identities, and
\eqref{inf:conditioning-final-inverse}.
All perturbations retain a factor $\tau$; only the final inverse
has a factor $\tau^{-1}$.
\end{proof}
